\documentclass[11pt,a4paper]{article}
\usepackage[utf8]{inputenc}
\usepackage{lmodern}
\usepackage{amsmath,amssymb,amsfonts}
\usepackage{booktabs}
\usepackage{geometry}
\usepackage{hyperref}
\usepackage{verbatim}

\geometry{margin=2.2cm}

\title{\textbf{Topological Emergence and Dimensional Projection of Time}}
\author{E. R. dos Santos}
\date{November 2026}

\begin{document}

\maketitle

\begin{abstract}
We present an exploratory phenomenological framework investigating the emergence of spatial scales and Standard Model mass hierarchies from discrete topological invariants and temporal boundary conditions. Starting from the Heisenberg energy-time uncertainty bound ($S_0 = 1/2$) and the Kolmogorov--Arnold--Moser (KAM) stability of the golden ratio ($\phi$) in non-linear pre-spatial dynamics, we examine how topological frustration in dense 3D ($z_{3d} = 12$) and 4D ($z_{4d} = 24$) lattice packings induces a discrete curvature defect $d^* = \frac{1}{2}\phi^{-1/3}$ and an Euler characteristic shift $\chi = 2 + \rho_v^* = 2.1$. Through an exponential dimensional projection of the Planck time ($t_{\mathrm{Planck}} \to \tau_c$), an infrared hadronic scale ($\mu_{\mathrm{IR}} \approx 234.97\,\mathrm{MeV}$) is obtained. We formulate an effective topological Lagrangian whose perturbative expansion yields both a Leading-Order (LO) tree-level spectrum and a Next-to-Leading-Order (NLO) spinor double-cover recoil correction across ten Standard Model observables, atomic benchmarks, and candidate beyond-Standard-Model (BSM) scales.
\end{abstract}

\section{Conceptual Postulates and Boundary Conditions}

Rather than assuming a pre-existing continuous spatial manifold, this investigation treats temporal evolution as the primary substrate from which spatial metric scales emerge via topological symmetry breaking. The boundary conditions of the model are structured around five working postulates:

\begin{enumerate}
    \item \textbf{Quantum Non-Nullity Bound:} A strictly static vacuum state ($\Delta E = 0$, $\Delta t = 0$) is precluded by the fundamental uncertainty relation $\Delta E \cdot \Delta t \ge \hbar / 2$. Consequently, the minimal primordial phase action is bounded below by $S_0 = 1/2$.
    \item \textbf{KAM Selection of the Golden Ratio:} In a non-linear pre-spatial regime governed by coupled temporal phase maps, rational frequency modes undergo resonant destabilization. By the Kolmogorov--Arnold--Moser (KAM) theorem, the invariant torus associated with the most irrational winding number---the golden ratio $\phi = (1+\sqrt{5})/2$---exhibits maximal resistance to chaotic breakdown.
    \item \textbf{Lattice Crystallization and Topological Frustration:} Upon phase coherence over $\phi$, discrete spatial coordination emerges in three-dimensional icosahedral/FCC packing ($z_{3d} = 12$) and four-dimensional $D_4$ (24-cell) lattice geometry ($z_{4d} = 24$). Because icosahedral symmetry cannot tile a spherical topology ($\chi_0 = 2$) without frustration, a discrete curvature defect and a topological vortex density arise:
    \begin{equation}
        d^* = \frac{1}{2}\phi^{-1/3} \approx 0.425900, \qquad \rho_v^* = \frac{1}{z_{3d} - \chi_0} = \frac{1}{10} = 0.1, \qquad \chi = \chi_0 + \rho_v^* = 2.1.
    \end{equation}
    \item \textbf{Dimensional Projection of Time:} The bare geometric gauge coupling $\alpha_{\mathrm{bare}}^{-1}$ governs an exponential scaling from the Planck time $t_{\mathrm{Planck}} = \sqrt{\hbar G / c^5}$ to a characteristic hadronic coherence time $\tau_c$, which projects onto an infrared energy scale $\mu_{\mathrm{IR}}$ via $\hbar c$:
    \begin{equation}
        \alpha_{\mathrm{bare}}^{-1} = \frac{2\pi \chi \sqrt{z_{4d}}}{d^*} \left(1 - \frac{d^*}{\sqrt{z_{4d}}} - \frac{\chi}{z_{3d}^2}\right), \qquad \tau_c = t_{\mathrm{Planck}} \exp\left(\frac{\alpha_{\mathrm{bare}}^{-1} - \chi / z_{3d}}{3}\right), \qquad \mu_{\mathrm{IR}} = \frac{\hbar c}{c \tau_c}.
    \end{equation}
    \item \textbf{Spinor Double-Cover Topology:} Spin-$1/2$ fermions are modeled as negative topological folds requiring a $4\pi$ double-cover rotation across the $2 z_{4d} = 48$ local spinorial states of the $D_4$ lattice.
\end{enumerate}

\section{Unified Effective Topological Lagrangian}

To connect the discrete lattice invariants with continuum field equations, we consider an effective action $S = \int d^4x \sqrt{-g_{\mathrm{eff}}}\,\mathcal{L}_{\mathrm{Total}}$ composed of four coupled sectors:
\begin{equation}
    \mathcal{L}_{\mathrm{Total}} = \mathcal{L}_{\mathrm{Lattice}}(\Phi, d^*) + \mathcal{L}_{\mathrm{Gauge}}(F_{\mu\nu}, \chi) + \mathcal{L}_{\mathrm{Spinor}}(\psi, \Phi, d^*) + \mathcal{L}_{\mathrm{Topological}}(a, G\tilde{G}).
\end{equation}

\subsection{Lattice Symmetry Breaking and Gauge-Hadronic Sectors}
The scalar lattice order parameter $\Phi$ acquires a vacuum expectation value $v_{\mathrm{Higgs}} = f_\pi \alpha^{-1}(0) z_{4d} \sqrt{\chi} \eta^{3/2}$, where $f_\pi = \mu_{\mathrm{IR}}/\sqrt{z_{4d}}$ and $\eta = (1 - d^*)^{-1/z_{3d}}$:
\begin{equation}
    \mathcal{L}_{\mathrm{Lattice}} = \frac{1}{2}\left(1 - \frac{d^*}{\sqrt{z_{4d}}} - \frac{\chi}{z_{3d}^2}\right)(D_\mu \Phi)^\dagger (D^\mu \Phi) - \lambda_H \left(\Phi^\dagger \Phi - \frac{v_{\mathrm{Higgs}}^2}{2}\right)^2, \qquad \lambda_H = \frac{z_{3d}}{4 z_{4d}}\eta.
\end{equation}
In the gauge and chiral hadronic sector, the weak mixing angle is fixed by the 3D coordination number ($\sin^2\theta_W = \frac{3}{z_{3d}+1} = \frac{3}{13}$), while baryonic states are stabilized as topological Skyrme solitons with $e_{\mathrm{Skyrme}} = (d^*)^{-1/4}$ and $c_{\mathrm{sol}} = 3^5/10 = 24.3$:
\begin{equation}
    \mathcal{L}_{\mathrm{Gauge}} = -\frac{1}{4 g_{\mathrm{bare}}^2} F_{\mu\nu}^a F^{a\,\mu\nu} + \frac{f_\pi^2}{4}\mathrm{Tr}(\partial_\mu U \partial^\mu U^\dagger) + \frac{1}{32 e_{\mathrm{Skyrme}}^2}\mathrm{Tr}[U^\dagger \partial_\mu U, U^\dagger \partial_\nu U]^2.
\end{equation}

\subsection{Spinor Kinetic Recoil: Perturbative Origin of LO and NLO}
Propagation of a spin-$1/2$ topological fold $\psi$ across the $2 z_{4d}$ double-cover states induces a curvature modification in the Dirac kinetic operator:
\begin{equation}
    \mathcal{L}_{\mathrm{Spinor}} = \bar{\psi}\left[ i\gamma^\mu D_\mu \left(1 + \frac{d^*}{2 z_{4d}} + \mathcal{O}(d^{*2})\right) - y_f(\eta, \chi)\Phi \right]\psi.
\end{equation}
Truncating at zeroth order in $d^*/(2z_{4d})$ defines the \textbf{Leading-Order (LO)} spectrum. Retaining the first-order kinetic term and canonically rescaling $\psi \to \left(1 + \frac{d^*}{2 z_{4d}}\right)^{-1/2}\psi$ introduces an elastic lattice recoil factor $\left(1 - \frac{d^*}{2 z_{4d}}\right)$ on physical charged-lepton masses, defining the \textbf{Next-to-Leading-Order (NLO)} spectrum. One-loop QED vacuum polarization is evaluated self-consistently up to $\Lambda = m_\tau$ (solved via the Koide relation $Q = 2/3$), while the top quark mass follows from vacuum saturation with 3D projection, $m_t = \frac{v_{\mathrm{Higgs}}}{\sqrt{2}}\left(1 - \frac{d^*}{z_{3d}^2}\right)$, entering $M_W$ and $M_Z$ via standard $\Delta\rho$ radiative corrections.

\section{Numerical Results and Empirical Comparison}

Table~\ref{tab:sm_results} summarizes the numerical outputs for both LO and NLO evaluations alongside Particle Data Group (PDG) reference values. Theoretical uncertainties ($\sigma_{\mathrm{tot}}$) combine renormalization scale variation ($\mu \in [0.5, 2.0]\,m_\tau$), higher-order radiative truncation ($\alpha/\pi$), and the experimental uncertainty of the gravitational constant ($\Delta G / G = 2.2 \times 10^{-5}$).

\begin{table}[ht!]
\centering
\caption{Comparison of Leading-Order (LO) and Next-to-Leading-Order (NLO) model outputs against experimental PDG values.}
\label{tab:sm_results}
\begin{tabular}{lcccc}
\toprule
\textbf{Observable} & \textbf{LO (Tree-Level)} & \textbf{NLO ($\pm \sigma_{\mathrm{tot}}$)} & \textbf{PDG Reference} & \textbf{NLO Tension} \\
\midrule
$M_p$ (GeV) & $0.9416$ & $0.9416 \pm 0.0022$ & $0.9383$ & $1.5\sigma$ \\
$\alpha^{-1}(0)$ & $137.1069$ & $137.1050 \pm 0.2943$ & $137.0360$ & $0.2\sigma$ \\
$m_e$ (MeV) & $0.5182$ & $0.5136 \pm 0.0016$ & $0.5110$ & $2.2\sigma$ \\
$m_\mu$ (MeV) & $106.5782$ & $105.6326 \pm 0.2452$ & $105.6584$ & $0.1\sigma$ \\
$m_\tau$ (MeV) & $1793.2838$ & $1777.3743 \pm 4.1399$ & $1776.8600$ & $0.1\sigma$ \\
$v_{\mathrm{Higgs}}$ (GeV) & $245.1447$ & $245.1413 \pm 0.7752$ & $246.2197$ & $1.4\sigma$ \\
$m_t$ (GeV) & $172.8308$ & $172.8284 \pm 0.5465$ & $172.6900$ & $0.3\sigma$ \\
$M_W$ (GeV) & $80.6378$ & $80.6373 \pm 0.2286$ & $80.3770$ & $1.1\sigma$ \\
$M_Z$ (GeV) & $91.2262$ & $91.2256 \pm 0.2586$ & $91.1876$ & $0.1\sigma$ \\
$M_H$ (GeV) & $125.4396$ & $125.4379 \pm 0.3966$ & $125.2500$ & $0.5\sigma$ \\
\bottomrule
\end{tabular}
\end{table}

\subsection{Emergent Quantum Scales and Testable Benchmarks}
Without introducing external QED loop coefficients, the geometric recoil of the defect $d^*$ yields an electron anomalous magnetic moment $a_e = \frac{\alpha}{2\pi} - d^*\left(\frac{\alpha}{\pi}\right)^2 = 0.0011585$ ($0.10\%$ relative deviation from $0.0011597$), an emergent Bohr radius $a_0 = 0.5267\,\text{\AA}$ ($0.46\%$ vs.\ $0.5292\,\text{\AA}$), and a Rydberg energy $R_\infty = 13.6621\,\mathrm{eV}$ ($0.41\%$ vs.\ $13.6057\,\mathrm{eV}$). Additionally, the lattice framework identifies two candidate beyond-Standard-Model benchmarks: a heavy radial scalar state at $M_{H2} = v_{\mathrm{Higgs}}\sqrt{z_{3d}/2} \approx 600.48\,\mathrm{GeV}$ (within reach of HL-LHC searches) and a GUT-scale topological axion with $f_a \approx 1.52 \times 10^{16}\,\mathrm{GeV}$ and $m_a \approx 0.376\,\mathrm{neV}$ (within the target sensitivity band of DMRadio-GUT).

\section*{Transparency and Methodology Note}
The conceptual formulation, theoretical axioms, and architecture of this model were idealized by \textbf{E. R. dos Santos}. The logical implementation in Python, mathematical structuring, and code refactoring were developed with the assistance of Artificial Intelligence.

This work does not aim, in any way or at any time, to exhaust the subject, supersede similar research, or provide definitive answers of any kind under any aspect. All numerical results, equations, and computational routines are available for reading, auditing, testing, and modification under the terms of the \textbf{GNU General Public License v3.0 (GPLv3)}.

\appendix
\section{Reference Python Implementation}
\footnotesize
\verbatiminput{generated/topological_emergence_model.py}

\end{document}
